\subsection{初項依存型}

最後に,漸化式の形だけでなく,初項の選び方そのものが問題を左右する
例を見よう.
\[
  a_1=c,
  \quad
  a_{n+1}=\frac{a_n^2-1}{2}
\]
初項を記号 $c$ のままにすると,
\begin{align*}
  a_2&=\frac{c^2-1}{2},\\
  a_3&=\frac{c^4-2c^2-3}{8},\\
  a_4&=\frac{c^8-4c^6-2c^4+12c^2-55}{128}
\end{align*}
一般に $a_n$ は $c$ の次数 $2^{n-1}$ の多項式になる.初項を一般の
文字のまま扱おうとすると,項を一つ進めるたびに次数が倍になり,
三角関数型や基本対称式型のような固定された表現が見えてこない.

特に $c=2$ と固定すると,
\[
  a_1=2,
  \quad
  a_2=\frac{3}{2},
  \quad
  a_3=\frac{5}{8},
  \quad
  a_4=-\frac{39}{128}
\]
となる.しかし,この列は直前の例のように角度を2倍する形にも,
指数の差を保つ形にも,接線近似の形にもならない.

ここでいう「$c=2$ 以外では解けない」とは,他の初項の数列が存在しない
という意味ではない.また,定数列を探すだけなら,
\[
  c=1+\sqrt{2},
  \quad
  c=1-\sqrt{2}
\]
では $a_{n+1}=a_n=c$ となる.大切なのは,ある初項を選んだときだけ
特別な恒等式や置換が働くことがあり,初項を一般化するとその偶然の
構造が失われる,という事実である.

また,初項依存型の例は先に挙げた基本対称式型でも同様に起こせる事象である.
もちろん,一般の初項での解法は先に示したとおりであるが,より狭義の解き方ではこのようなものがある.
\[
  a_1=-1,\quad a_{n+1}=a_n^2-2
\]
のような場合であれば,第5項までが以下のようになる.
\[
  -1,-1,-1,-1,-1
\]
漸化式が表す数列が定数数列を成したのだ.
このような解はシンプルに不動点を求めることで作り出せる.
つまりは
\[
  \alpha=\alpha^2-2
\]
を$\alpha$について解くけばよい.
この二次方程式が成り立つような$\alpha$は$-1$以外に2が存在する.
実際に代入してみると常に$a_n=2$が成り立つことが分かる.

この例では,少なくとも本節で扱った初等的な道具からは一般項を得る
決まった方法が出てこない.それでも漸化式は数列を定めており,
\[
  a_n=f^{\,\circ(n-1)}(c),
  \quad
  f(x)=\frac{x^2-1}{2}
\]
と書くことはできる.これは解答の終わりではなく,「一般項を求める」
という目標自体を見直す入口である.付録では,閉じた式が得られない
場合にも,軌道,極限,近似,行列,母関数など別の言葉で数列を理解する
方法を考える.
